\documentclass[11pt,a4paper]{article}

% --- Mathematics (must precede unicode-math) ---
\usepackage{amsmath}
\usepackage{mathtools}

% --- Fonts (lualatex) ---
\usepackage{fontspec}
\usepackage{unicode-math}
\setmathfont{KpMath-Regular.otf}
\setmathfont{KpMath-Bold.otf}[version=bold]

% --- Microtypography (after all font setup) ---
\usepackage{microtype}

% --- Colors & page layout ---
\usepackage[dvipsnames]{xcolor}
\usepackage[margin=25mm, top=30mm, bottom=30mm]{geometry}

% --- Headers & footers ---
\usepackage{fancyhdr}
\pagestyle{fancy}
\fancyhf{}
\fancyhead[L]{\small\itshape Multipolar Effective Sources and Mie Theory}
\fancyhead[R]{\small\itshape Nestor (2026)}
\fancyfoot[C]{\thepage}
\renewcommand{\headrulewidth}{0.4pt}
\renewcommand{\footrulewidth}{0pt}

% --- Section numbering ---
\setcounter{secnumdepth}{3}

% --- Tables ---
\usepackage{booktabs}
\usepackage{array}

% --- Captions ---
\usepackage[font=small, labelfont=bf, labelsep=period]{caption}

% --- Spacing ---
\usepackage{setspace}
\setstretch{1.15}
\setlength{\parskip}{0.4em}

% --- Hyperlinks (last before macros) ---
\usepackage[colorlinks=true, linkcolor=blue!60!black,
            citecolor=blue!60!black, urlcolor=blue!60!black]{hyperref}

% --- Macros (reused from cubic_tmatrix.tex) ---
\newcommand{\bG}{\mathbf{G}}
\newcommand{\bI}{\mathbf{I}}
\newcommand{\bx}{\mathbf{x}}
\newcommand{\beR}{\hat{\mathbf{e}}_R}
\newcommand{\dd}{\mathrm{d}}
\newcommand{\order}[1]{\mathcal{O}\!\left(#1\right)}

\newcommand{\bT}{\mathbf{T}}
\newcommand{\bE}{\mathbf{E}}
\newcommand{\bS}{\mathbf{S}}
\newcommand{\Dlam}{\Delta\lambda}
\newcommand{\Dmu}{\Delta\mu}
\newcommand{\Drho}{\Delta\rho}
\newcommand{\Dkappa}{\Delta\kappa}
\newcommand{\PV}{\operatorname{P.V.}}
\newcommand{\re}{\operatorname{Re}}
\newcommand{\im}{\operatorname{Im}}
\newcommand{\tr}{\operatorname{tr}}

% --- New macros for this document ---
\newcommand{\bu}{\mathbf{u}}
\newcommand{\bn}{\hat{\mathbf{n}}}
\newcommand{\bk}{\mathbf{k}}
\newcommand{\beps}{\boldsymbol{\varepsilon}}
\newcommand{\bsig}{\boldsymbol{\sigma}}

\begin{document}

% --- Title block ---
\thispagestyle{plain}
\begin{center}
\vspace*{1em}
{\LARGE\bfseries Multipolar Effective Sources\\[0.3em]
and Elastic Mie Theory}\\[1.8em]
{\large T.\,M.\ Nestor}\\[0.5em]
{\itshape Research School of Earth Sciences, Australian National University}\\[2em]
\end{center}

% =============================================================================
\section{Introduction}
\label{sec:intro}
% =============================================================================

Elastic scattering from a sphere embedded in a homogeneous isotropic
medium admits two complementary descriptions.

The first is \emph{elastic Mie theory}: a partial wave expansion in
spherical harmonics that exactly satisfies the boundary conditions at the
sphere surface~\cite{YingTruell1956,PaoMow1973,EinspruchWitterholtTruell1960}.
The solution is expressed as a set of scattering coefficients $a_n$,
$b_n$, $c_n$ for each angular order~$n$, determined by a $4 \times 4$
(P-SV) or $2 \times 2$ (SH) linear system per order.  This is exact at
all frequencies and contrasts but provides limited physical insight into
the scattering mechanism.

The second is the \emph{Eshelby T-matrix} formulation: the scattered
field is generated by volume-averaged \emph{effective sources} inside
the scatterer---an equivalent body force density (from the density
contrast) and an equivalent stress perturbation (from the elastic
modulus contrast)~\cite{Eshelby1957,GubernatisDomanyKrumhansl1977,
Waterman1969}.  Self-consistent amplification factors account for the
interaction between the interior field and the scatterer's own
radiation, yielding effective contrasts $\Drho^*$, $\Dkappa^*$, and
$\Dmu^*$.

The central result of this document is the \emph{exact mapping} between
these two descriptions: each effective source type radiates with a
distinct multipolar pattern, and the leading Mie partial waves
$n = 0, 1, 2$ correspond one-to-one to the three independent effective
contrasts.  This mapping is validated numerically by the companion
Python module \texttt{sphere\_scattering.py}.

% =============================================================================
\section{Elastic Mie Theory}
\label{sec:mie}
% =============================================================================

\subsection{Helmholtz decomposition}

Consider a sphere of radius~$a$ with interior properties
$(\rho_1, \lambda_1, \mu_1)$ embedded in a background medium
$(\rho_0, \lambda_0, \mu_0)$ with P-wave velocity
$\alpha_0 = \sqrt{(\lambda_0 + 2\mu_0)/\rho_0}$ and S-wave velocity
$\beta_0 = \sqrt{\mu_0/\rho_0}$.

The displacement field is decomposed via Helmholtz potentials:
\begin{equation}
  \bu = \nabla\varphi + \nabla \times \nabla \times (\hat{\mathbf{r}}\,r\,\psi)
    + \nabla \times (\hat{\mathbf{r}}\,r\,\chi),
  \label{eq:helmholtz}
\end{equation}
where $\varphi$ generates the P-wave (compressional) part, $\psi$
generates the SV-wave part, and $\chi$ generates the SH-wave part.
Each potential satisfies a scalar Helmholtz equation:
\begin{equation}
  (\nabla^2 + k_P^2)\,\varphi = 0, \qquad
  (\nabla^2 + k_S^2)\,\psi = 0, \qquad
  (\nabla^2 + k_S^2)\,\chi = 0,
  \label{eq:helmholtz_eq}
\end{equation}
with $k_P = \omega/\alpha$ and $k_S = \omega/\beta$ for the
appropriate medium (interior or exterior).

\subsection{Partial wave expansion}

For an incident P-wave $\bu_\text{inc} = \hat{\mathbf{z}}\,
e^{ik_P^0 z}$ propagating along the $z$-axis, the scalar potential
is $\varphi_\text{inc} = e^{ik_P^0 z}/(ik_P^0)$.  The plane wave
expansion in Legendre polynomials reads
\begin{equation}
  e^{ik_P^0 r\cos\theta}
  = \sum_{n=0}^{\infty} (2n+1)\,i^n\,j_n(k_P^0 r)\,P_n(\cos\theta),
  \label{eq:plane_wave}
\end{equation}
where $j_n$ is the spherical Bessel function and $P_n$ is the Legendre
polynomial of degree~$n$.

The incident potential is therefore
\begin{equation}
  \varphi_\text{inc}
  = \sum_{n=0}^{\infty} \frac{(2n+1)\,i^n}{ik_P^0}\,
    j_n(k_P^0 r)\,P_n(\cos\theta).
  \label{eq:phi_inc}
\end{equation}

The scattered and interior fields are expanded similarly:
\begin{alignat}{2}
  \varphi_\text{scat} &= \sum_n A_n\,h_n^{(1)}(k_P^0 r)\,P_n,
  &\qquad
  \psi_\text{scat} &= \sum_n B_n\,h_n^{(1)}(k_S^0 r)\,P_n,
  \label{eq:scat_potentials} \\[4pt]
  \varphi_\text{int}  &= \sum_n C_n\,j_n(k_P^1 r)\,P_n,
  &\qquad
  \psi_\text{int}  &= \sum_n D_n\,j_n(k_S^1 r)\,P_n,
  \label{eq:int_potentials}
\end{alignat}
where $h_n^{(1)}$ is the spherical Hankel function of the first kind
(outgoing waves), and the superscripts $0$, $1$ denote the exterior
and interior media respectively.

\subsection{Boundary conditions}

At $r = a$, continuity of displacement and traction requires:
\begin{equation}
  \begin{cases}
    u_r^\text{scat} + u_r^\text{inc} = u_r^\text{int}, \\[3pt]
    u_\theta^\text{scat} + u_\theta^\text{inc} = u_\theta^\text{int}, \\[3pt]
    \sigma_{rr}^\text{scat} + \sigma_{rr}^\text{inc}
      = \sigma_{rr}^\text{int}, \\[3pt]
    \sigma_{r\theta}^\text{scat} + \sigma_{r\theta}^\text{inc}
      = \sigma_{r\theta}^\text{int}.
  \end{cases}
  \label{eq:bc}
\end{equation}
The displacement and stress components from the P-wave potential
$\varphi = z_n(kr)\,P_n(\cos\theta)$ are:
\begin{align}
  u_r &= k\,z_n'(kr)\,P_n, &
  u_\theta &= \frac{z_n(kr)}{r}\,\frac{\dd P_n}{\dd\theta},
  \label{eq:ur_P}
\end{align}
\begin{align}
  \sigma_{rr} &= \Bigl[
    -(\lambda + 2\mu)\,k^2 z_n
    - \frac{4\mu k}{r}\,z_n'
    + \frac{2\mu n(n{+}1)}{r^2}\,z_n
  \Bigr]\,P_n,
  \label{eq:srr_P} \\[4pt]
  \sigma_{r\theta} &= \frac{2\mu}{r}
    \Bigl[k\,z_n' - \frac{z_n}{r}\Bigr]\,
    \frac{\dd P_n}{\dd\theta}.
  \label{eq:srt_P}
\end{align}
The SV-wave potential $\psi = z_n(kr)\,P_n$ contributes:
\begin{align}
  u_r &= \frac{n(n{+}1)\,z_n}{r}\,P_n, &
  u_\theta &= \frac{z_n + kr\,z_n'}{r}\,
    \frac{\dd P_n}{\dd\theta},
  \label{eq:ur_S}
\end{align}
\begin{align}
  \sigma_{rr} &= \frac{2\mu n(n{+}1)}{r}
    \Bigl[k\,\frac{z_n'}{r} - \frac{z_n}{r^2}\Bigr]\,P_n,
  \label{eq:srr_S} \\[4pt]
  \sigma_{r\theta} &= \frac{\mu}{r^2}
    \Bigl[(2n^2{+}2n{-}2{-}k^2r^2)\,z_n
    - 2kr\,z_n'\Bigr]\,
    \frac{\dd P_n}{\dd\theta}.
  \label{eq:srt_S}
\end{align}

\subsection{Scattering matrix}
\label{sec:mie_matrix}

Substituting the expansions~\eqref{eq:scat_potentials}--\eqref{eq:int_potentials}
into the boundary conditions~\eqref{eq:bc}, each angular order~$n$ decouples
into a $4 \times 4$ linear system for the P-SV coefficients:
\begin{equation}
  \begin{pmatrix}
    u_r^{P,h} & u_r^{S,h} & -u_r^{P,j} & -u_r^{S,j} \\[3pt]
    u_\theta^{P,h} & u_\theta^{S,h} & -u_\theta^{P,j} & -u_\theta^{S,j} \\[3pt]
    \sigma_{rr}^{P,h} & \sigma_{rr}^{S,h} & -\sigma_{rr}^{P,j} & -\sigma_{rr}^{S,j} \\[3pt]
    \sigma_{r\theta}^{P,h} & \sigma_{r\theta}^{S,h} & -\sigma_{r\theta}^{P,j} & -\sigma_{r\theta}^{S,j}
  \end{pmatrix}
  \begin{pmatrix} A_n \\ B_n \\ C_n \\ D_n \end{pmatrix}
  = -\begin{pmatrix} u_r^\text{inc} \\ u_\theta^\text{inc} \\
      \sigma_{rr}^\text{inc} \\ \sigma_{r\theta}^\text{inc} \end{pmatrix},
  \label{eq:mie_4x4}
\end{equation}
where the superscripts $(P,h)$, $(S,h)$ denote the scattered P and
S fields (outgoing Hankel functions, exterior material parameters) and
$(P,j)$, $(S,j)$ denote the interior fields (regular Bessel functions,
interior material parameters), all evaluated at $r = a$.

The monopole $n = 0$ is purely compressional (no shear contribution)
and reduces to a $2 \times 2$ system.  The SH modes decouple entirely
and are governed by a separate $2 \times 2$ system for each
$n \geq 1$~\cite{PaoMow1973}.

\subsection{Far-field scattering amplitude}

The P-wave scattering coefficients $a_n \equiv A_n$ define the
far-field amplitude via:
\begin{equation}
  f_P(\theta) = \sum_{n=0}^{\infty} a_n\,(-i)^n\,P_n(\cos\theta),
  \qquad
  \bu_P^\text{scat} \sim f_P(\theta)\,\hat{\mathbf{r}}\,
    \frac{e^{ik_P r}}{r}.
  \label{eq:farfield}
\end{equation}
Note: the stored coefficients $a_n$ include a factor $(-1)^n$ relative
to the standard convention, so that the scattering is forward-peaked
in our coordinate system (axis~0 = propagation
direction)~\cite{KorneevJohnson1993}.

% =============================================================================
\section{Eshelby T-Matrix for a Sphere}
\label{sec:eshelby}
% =============================================================================

\subsection{Equivalent body sources}

The Lippmann-Schwinger equation recasts the scattering problem as
radiation from equivalent sources distributed throughout the scatterer
volume~$V$~\cite{GubernatisDomanyKrumhansl1977}.  For an isotropic
sphere with contrasts $\Drho = \rho_1 - \rho_0$,
$\Dlam = \lambda_1 - \lambda_0$, $\Dmu = \mu_1 - \mu_0$:

\paragraph{Inertial source.}
The density contrast generates an equivalent body force:
\begin{equation}
  f_i(\bx) = -\omega^2\,\Drho\,u_i(\bx),
  \label{eq:force_source}
\end{equation}
proportional to the total displacement field inside the scatterer.

\paragraph{Elastic source.}
The modulus contrast generates an equivalent stress perturbation:
\begin{equation}
  \Delta\sigma_{ij}(\bx) = \Delta C_{ijkl}\,\varepsilon_{kl}(\bx),
  \label{eq:stress_source}
\end{equation}
where $\Delta C_{ijkl} = \Dlam\,\delta_{ij}\delta_{kl}
+ \Dmu\,(\delta_{ik}\delta_{jl} + \delta_{il}\delta_{jk})$ is the
isotropic stiffness contrast tensor and $\varepsilon_{kl}$ is the
total strain inside the scatterer.

\subsection{Volume-averaged effective sources}

In the Rayleigh regime ($ka \ll 1$), the internal fields are
approximately uniform and can be replaced by their volume averages.
The scattered field is then generated by:
\begin{itemize}
  \item a \emph{force monopole}
    $F_i = -\omega^2 V \Drho^*\,u_i^\text{inc}$
    (uniform body force $\times$ volume), and
  \item a \emph{stress dipole}
    $\Delta\sigma^*_{ij} = V\,\Delta C^*_{ijkl}\,
    \varepsilon_{kl}^\text{inc}$
    (uniform stress perturbation $\times$ volume),
\end{itemize}
where the starred quantities are the \emph{self-consistent effective
contrasts} that account for the modification of the internal field by
the scatterer's own radiation.

\subsection{Decomposition into isotropic and deviatoric parts}

The effective stress source decomposes as
\begin{equation}
  \Delta\sigma^*_{ij}
  = \Dkappa^*\,\theta\,\delta_{ij}
  + 2\Dmu^*\,\varepsilon'_{ij},
  \label{eq:stress_decomp}
\end{equation}
where $\theta = \tr(\beps) = \varepsilon_{kk}$ is the dilatation,
$\varepsilon'_{ij} = \varepsilon_{ij} - \tfrac{1}{3}\theta\,\delta_{ij}$
is the deviatoric strain, and
\begin{equation}
  \Dkappa^* = \Dlam^* + \tfrac{2}{3}\Dmu^*
  \label{eq:Dkappa_def}
\end{equation}
is the effective bulk modulus contrast.

\subsection{Eshelby depolarisation tensor}

The self-consistent amplification factors are determined by the
volume integral of the Green's tensor second derivative over the
sphere~\cite{Eshelby1957,Mura1987}:
\begin{equation}
  \int_V \frac{\partial^2 G_{ij}}{\partial x_k\,\partial x_l}\,\dd V
  = A\,\delta_{ij}\delta_{kl}
  + B\,(\delta_{ik}\delta_{jl} + \delta_{il}\delta_{jk}).
  \label{eq:IVsphere}
\end{equation}
For a sphere, rotational symmetry eliminates the cubic anisotropy
term ($C = 0$) present for non-spherical shapes.

The integral~\eqref{eq:IVsphere} has a critical subtlety: the
$1/r$ singularity of the Green's tensor produces a
\emph{delta function} under second differentiation.  This Eshelby
contribution dominates the integral and must be treated analytically.
For an isotropic sphere in an isotropic
matrix~\cite{Eshelby1957,Mura1987}:
\begin{equation}
  \boxed{
  A_E = -\frac{9 - 10\nu}{30\mu_0(1 - \nu)},
  \qquad
  B_E = \frac{1}{30\mu_0(1 - \nu)},
  }
  \label{eq:eshelby_AB}
\end{equation}
where $\nu = \lambda_0 / [2(\lambda_0 + \mu_0)]$ is the background
Poisson ratio.

The smooth (dynamic) correction from the oscillating exponentials
$e^{i\omega r/\alpha}$, $e^{i\omega r/\beta}$ is
$\order{(ka)^2}$ relative to the Eshelby part and is negligible
in the Rayleigh regime.

\subsection{Self-consistent amplification factors}

The T-matrix coupling coefficients are
\begin{equation}
  T_1 = A\,\Dlam + 2B\,\Dmu, \qquad
  T_2 = B\,\Dlam + (A + B)\,\Dmu + B\,\Dmu,
\end{equation}
and the amplification factors follow from:
\begin{align}
  A_u &= \frac{1}{1 - \omega^2\,\Drho\,\Gamma_0},
  \label{eq:amp_u} \\[4pt]
  A_\theta &= \frac{1}{1 - 3T_1 - 2T_2},
  \label{eq:amp_theta} \\[4pt]
  A_\text{dev} &= \frac{1}{1 - 2T_2},
  \label{eq:amp_dev}
\end{align}
where $\Gamma_0 = \int_V G_{ii}\,\dd V / 3$ is the scalar volume
integral of the Green's tensor.  Here $A_u$ amplifies the
displacement (governing the effective density), $A_\theta$ amplifies
the dilatation (governing the effective bulk modulus), and
$A_\text{dev}$ amplifies the deviatoric strain (governing the
effective shear modulus).

The effective contrasts are then:
\begin{align}
  \Drho^*   &= \Drho \cdot A_u,
  \label{eq:Drho_star} \\[4pt]
  \Dkappa^* &= (\Dlam + \tfrac{2}{3}\Dmu)\,A_\theta
             = \Dkappa\,A_\theta,
  \label{eq:Dkappa_star} \\[4pt]
  \Dmu^*    &= \Dmu \cdot A_\text{dev}.
  \label{eq:Dmu_star}
\end{align}

% =============================================================================
\section{Multipolar Radiation from Effective Sources}
\label{sec:multipolar}
% =============================================================================

This section establishes the central connection: each type of effective
source radiates with a distinct multipolar angular pattern, and the
leading Mie partial waves map one-to-one to the three effective
contrasts.

\subsection{Rayleigh scattering formula}

The far-field P-wave displacement scattered by the effective sources
in a sphere of volume $V = \tfrac{4}{3}\pi a^3$
is~\cite{AkiRichards2002,GubernatisDomanyKrumhansl1977}:
\begin{equation}
  f_P(\theta) = -\frac{V}{4\pi\rho_0\alpha_0^2}\,\Bigl[
    k_P^2\,\Dkappa^*
    + \omega^2\,\Drho^*\,\cos\theta
    + 2k_P^2\,\Dmu^*\,\cos^2\!\theta
  \Bigr].
  \label{eq:rayleigh_f}
\end{equation}
Using the Legendre polynomial identities $P_0 = 1$,
$P_1(\cos\theta) = \cos\theta$, and
$\cos^2\!\theta = \tfrac{1}{3}(2P_2 + 1)$, this becomes:
\begin{equation}
  \boxed{
  f_P(\theta) = -C_P\,\Bigl[
    k_P^2\,\Dkappa^*\,P_0
    + \omega^2\,\Drho^*\,P_1
    + \tfrac{4}{3}k_P^2\,\Dmu^*\,P_2
  \Bigr],
  }
  \label{eq:rayleigh_legendre}
\end{equation}
where $C_P = V/(4\pi\rho_0\alpha_0^2)$ and the isotropic part of
the $\Dmu^*$ term (the $\tfrac{2}{3}k_P^2\Dmu^*$ from
$\cos^2\!\theta \to P_0/3$) has been absorbed into the effective
$\Dkappa^*$ via~\eqref{eq:Dkappa_def}.

\subsection{Monopole radiation: isotropic stress \texorpdfstring{($n = 0$)}{(n=0)}}

The isotropic stress source $\Dkappa^*\,\theta\,\delta_{ij}$
produces a spherically symmetric pressure field---pure
\emph{monopole} radiation.  This couples exclusively to the $n = 0$
Legendre mode.  Physically, a uniform dilatation inside the sphere
acts as a breathing source, expanding and contracting isotropically.

Comparing the $P_0$ coefficient of~\eqref{eq:rayleigh_legendre}
with the Mie expansion~\eqref{eq:farfield}:
\begin{equation}
  a_0 = -C_P\,k_P^2\,\Dkappa^*.
  \label{eq:a0_match}
\end{equation}

\subsection{Dipole radiation: body force \texorpdfstring{($n = 1$)}{(n=1)}}

The uniform body force $f_i = -\omega^2\,\Drho^*\,u_i^\text{inc}$
produces a $\cos\theta$ radiation pattern---pure \emph{dipole}
radiation.  A uniform force density applied to a sphere translates
it bodily; the radiation is that of an oscillating point force
(elastodynamic analogue of an electric dipole).  This couples
exclusively to the $n = 1$ Legendre mode.

Matching the $P_1$ coefficient:
\begin{equation}
  -i\,a_1 = -C_P\,\omega^2\,\Drho^*.
  \label{eq:a1_match}
\end{equation}
The factor of $(-i)$ arises from the $(-i)^n$ phase in the Mie
far-field expansion~\eqref{eq:farfield} at $n = 1$.

\subsection{Quadrupole radiation: deviatoric stress \texorpdfstring{($n = 2$)}{(n=2)}}

The deviatoric stress source $2\Dmu^*\,\varepsilon'_{ij}$ produces a
$\cos^2\!\theta$ radiation pattern---pure \emph{quadrupole} radiation.
A deviatoric strain applied uniformly to a sphere deforms it into an
oblate/prolate ellipsoid; the radiation has the angular dependence of
a force couple without moment (elastodynamic analogue of an electric
quadrupole).  This couples exclusively to the $n = 2$ Legendre mode.

Matching the $P_2$ coefficient:
\begin{equation}
  -a_2 = -C_P\,\tfrac{4}{3}k_P^2\,\Dmu^*.
  \label{eq:a2_match}
\end{equation}
The factor of $(-1) = (-i)^2$ comes from the $n = 2$ phase.

\subsection{Extraction formulas}

Inverting~\eqref{eq:a0_match}--\eqref{eq:a2_match} gives the
\emph{exact mapping} from Mie coefficients to effective contrasts:
\begin{equation}
  \boxed{
  \Dkappa^* = -\frac{a_0}{C_P\,k_P^2},
  \qquad
  \Drho^* = \frac{i\,a_1}{C_P\,\omega^2},
  \qquad
  \Dmu^* = \frac{3\,a_2}{4\,C_P\,k_P^2}.
  }
  \label{eq:extraction}
\end{equation}
The effective Lam\'e contrast follows from $\Dlam^* = \Dkappa^* -
\tfrac{2}{3}\Dmu^*$.

An independent density estimate comes from the S-wave dipole
coefficient $b_1$ via $\Drho^*_S = i\,b_1 / (C_S\,\omega^2)$
with $C_S = V/(4\pi\rho_0\beta_0^2)$.

\subsection{Why higher multipoles are absent}

The T-matrix with volume-averaged effective sources generates
\emph{exactly three} independent radiation patterns: monopole, dipole,
and quadrupole.  This is a direct consequence of the source structure:
\begin{itemize}
  \item A \emph{uniform} body force (rank-1 tensor, 3 components)
    radiates as a dipole ($n = 1$ only).
  \item A \emph{uniform isotropic} stress (rank-0, scalar $\Dkappa^*$)
    radiates as a monopole ($n = 0$ only).
  \item A \emph{uniform deviatoric} stress (rank-2 traceless symmetric
    tensor, 5 components) radiates as a quadrupole ($n = 2$ only).
\end{itemize}
Higher multipoles ($n \geq 3$) arise from \emph{spatial variation}
of the internal field, which the volume-averaging approximation discards.
In the Rayleigh limit ($ka \to 0$), the internal field is genuinely
uniform, so $a_n = \order{(ka)^{2n+1}}$ for $n \geq 3$ and the
T-matrix result is exact.

% =============================================================================
\section{Convergence and Limitations}
\label{sec:convergence}
% =============================================================================

\subsection{Internal field non-uniformity at finite \texorpdfstring{$ka$}{ka}}

As $ka$ increases, the incident wavelength becomes comparable to
the sphere diameter and the internal field develops spatial structure
(gradients, standing wave patterns).  The volume-averaged effective
sources then fail to capture the higher multipoles excited by these
non-uniform internal fields.  The Mie solution, which enforces
boundary conditions exactly, automatically includes all multipole
orders.

The practical consequence is that the effective contrasts extracted
via~\eqref{eq:extraction} from Mie coefficients remain meaningful
(they correctly give the $n = 0, 1, 2$ contributions), but they no
longer fully characterise the scattered field because the $n \geq 3$
terms become significant.

\subsection{Role of contrast strength}

At weak contrast ($|\Delta C/C_0| \ll 1$), the Born approximation
holds: the internal field equals the incident field, and the
amplification factors $A_u \approx A_\theta \approx A_\text{dev}
\approx 1$.  The T-matrix effective contrasts then equal the bare
contrasts, and the extraction formulas~\eqref{eq:extraction} give
\emph{exact} agreement between the two routes.

At finite contrast, the self-consistent amplification factors
modify the effective contrasts.  The Eshelby route (volume integral
of the Green's tensor second derivative) and the Mie route
(boundary matching) produce slightly different amplification
factors because:
\begin{enumerate}
  \item The Eshelby approach assumes uniform interior fields and
    corrects \emph{globally} via volume-averaged self-interaction.
  \item The Mie approach satisfies boundary conditions exactly but
    does not assume interior uniformity.
\end{enumerate}
In the Rayleigh limit, both give the same result.  At finite $ka$,
the Eshelby route inherits a systematic bias of $\order{(ka)^2}$
from the uniform-field assumption.

\subsection{Positive vs.\ negative contrasts}

Positive contrasts (stiffer, denser inclusion) are generally better
behaved than negative contrasts of the same magnitude.  Negative
contrasts can produce resonances (the interior wavelength shortens,
leading to standing waves inside the sphere at relatively low $ka$)
which excite high-order multipoles and break the volume-averaging
assumption earlier.

\subsection{Foldy-Lax voxelization as a bridge}

The Foldy-Lax voxelization approach (see companion document
\emph{Coupled Foldy-Lax T-Matrix}) recovers higher multipoles
($n \geq 3$) through inter-cell multiple scattering within the
voxelized sphere.  Each sub-cell is in the deep Rayleigh regime
and correctly generates only $n = 0, 1, 2$ radiation, but the
mutual interaction between many sub-cells creates the complex
internal field structure needed to reproduce $n \geq 3$ partial
waves~\cite{Foldy1945,Lax1951,Lax1952}.

This provides a systematic convergence path:
\begin{itemize}
  \item $n_\text{sub} = 1$: single cell, Rayleigh T-matrix,
    $n = 0, 1, 2$ only.
  \item $n_\text{sub} = 3$: $N \approx 19$ cells (inside sphere),
    captures $n \leq 4$--$5$ with moderate accuracy.
  \item $n_\text{sub} = 5$--$7$: $N \approx 81$--$179$ cells,
    converges to the Mie solution at moderate $ka$.
\end{itemize}

% =============================================================================
\section{Summary}
\label{sec:summary}
% =============================================================================

Table~\ref{tab:summary} collects the complete correspondence between
multipole order, physical mechanism, effective contrast, and extraction
formula.

\begin{table}[ht]
\centering
\caption{Mapping between Mie multipole orders, effective sources,
  and extraction formulas.  Here
  $C_P = V/(4\pi\rho_0\alpha_0^2)$,
  $C_S = V/(4\pi\rho_0\beta_0^2)$,
  $V = \tfrac{4}{3}\pi a^3$, and $k_P = \omega/\alpha_0$.}
\label{tab:summary}
\begin{tabular}{@{}clllll@{}}
\toprule
$n$ & Name & Source type & Contrast & Extraction & Scaling \\
\midrule
0 & Monopole & Isotropic stress
  & $\Dkappa^*$ & $-a_0/(C_P k_P^2)$
  & $(ka)^3$ \\[4pt]
1 & Dipole & Body force (inertial)
  & $\Drho^*$ & $i\,a_1/(C_P \omega^2)$
  & $(ka)^3$ \\[4pt]
2 & Quadrupole & Deviatoric stress
  & $\Dmu^*$ & $3\,a_2/(4\,C_P k_P^2)$
  & $(ka)^5$ \\[4pt]
$\geq 3$ & Higher & Interior field gradients
  & --- & (not captured by T-matrix)
  & $(ka)^{2n+1}$ \\
\bottomrule
\end{tabular}
\end{table}

The Mie scattering coefficients $a_n$ scale as $(ka)^{2n+1}$ in the
Rayleigh limit~\cite{YingTruell1956}, confirming that the monopole
and dipole terms ($n = 0, 1$) dominate at long wavelengths.  The
quadrupole ($n = 2$) enters at $\order{(ka)^5}$, while $n \geq 3$
contributions are $\order{(ka)^7}$ and below.

The effective contrast extraction~\eqref{eq:extraction} is
\emph{exact} in the Rayleigh limit ($ka \to 0$) for all contrast
magnitudes, and remains a useful approximation into the transition
regime ($ka \sim 0.3$--$0.5$).  Beyond this, the Foldy-Lax
voxelization provides a systematic extension that recovers the full
Mie solution through emergent multipolar scattering from inter-cell
interactions.

\bibliographystyle{plain}
\bibliography{multipolar_mie}

\end{document}
